% !PRINTED 697
\weberpage{697}
%<<CONT
\weq{5}{
\begin{aligned}
(O, B) &= (r_{s+1} - 2) + (r_{s+2} - 2) + \cdots + (r_n - 2), \\
p &= (r_{s+1} - 1) + (r_{s+2} - 1) + \cdots + (r_n - 1), \\
(O, B) &= p - n + s,
\end{aligned}
}
and by (2):
\[
(O, A) - (O, B) = n - p + 1.
\]

Now by \S\,192, (14) $(O, B) = (A, W)$, and we have
the \gesperrt{Riemann-Roch} theorem for classes of the first kind:

4. \gesperrt{The dimension of a proper class $A$
of the first kind is}
\weq{6}{(O, A) = n - p + 1 + (A, W).}

This formula contains (3) within itself, because for polygons of the
second kind $(A, W) = 0$.

Since the dimension of a proper class is at least equal
to 2, it follows for a class of the second kind that
\[
n \geqq p + 1,
\]
and herein is contained the proof of a theorem of \gesperrt{Riemann}:

5. \gesperrt{A function whose order is less than
$(p + 1)$ is of the first kind.}

It follows further from this that the class $W$ of the polygons
of the first kind is a proper one. For suppose $W$ had
the divisor $\mathfrak{M}$, and let $W = \mathfrak{M} A$; then $(O, W) = (O, A) = p$.
If one takes $\mathfrak{A}$ in $A$ relatively prime to $\mathfrak{M}$, then in $W$ there is only
one polygon divisible by $\mathfrak{A}$, namely $\mathfrak{A}\mathfrak{M}$, and consequently
$(A, W) = 1$. Accordingly formula (6) gives:
\[
n = 2p - 2,
\]
hence equal to the order of $W$. The order of $A$ is therefore
just as great as that of $W$, and consequently $\mathfrak{M} = \mathfrak{O}$.

Theorem 4. holds unchanged for improper classes as well.
In order to prove this, let $A$ first be of the first kind
and of divisor $\mathfrak{M}$:
\weq{7}{
\begin{aligned}
A &= \mathfrak{M} A', \\
AB &= W, \\
A'B' &= W, \\
B' &= \mathfrak{M} B, \quad (\S\,192,\,8.) \\
(A', W) &= (O, B').
\end{aligned}
}

Let the orders of $\mathfrak{M}, A, B$ be $m, a, b$, so that $a + b$ is the
order of $W$, i.e.
\weq{8}{a + b = 2(p - 1).}
